\documentclass[]{beamer}
\usepackage{pgfpages}

% \setbeameroption{show notes on second screen}

% Include any extra LaTeX packages required
\usepackage[
  backend=biber,
  style=alphabetic,
  citestyle=authoryear
]{biblatex}

\addbibresource{./library.bib}

%% Packages
  \usepackage{listings}
  \usepackage{multicol}
  \usepackage{booktabs}
  \usepackage{xcolor}
  \usepackage{times}
  \usepackage{tikz}
  \usepackage{amsmath}
  \usepackage{verbatim}
  \usetikzlibrary{arrows.meta,shapes}

\usetheme{metropolis}

%% Settings and new commands
  \def\mAlertSpace{\vspace{0.5em}}
  \def\mOrangeItem{\item[\textcolor{mLightBrown}{\textbullet}]}
  % 36616a
  \definecolor{feldgrau}{rgb}{0.3, 0.36, 0.33}
  \newcommand{\mSlideTitle}{{{\color{gray}\secname}} \# \subsecname \hfill {\scriptsize$|$ UNED}}
  \newcommand{\mShortTitle}{\secname \hfill {\scriptsize$|$ UNED}}

  \tikzstyle{every picture}+=[remember picture]
  \everymath{\displaystyle}
  \definecolor{eclipseStrings}{RGB}{42,0.0,255}
\definecolor{eclipseKeywords}{RGB}{127,0,85}
\colorlet{numb}{magenta!60!black}

\lstdefinelanguage{json}{%
    basicstyle=\normalfont\ttfamily,
    commentstyle=\color{eclipseStrings}, % style of comment
    stringstyle=\color{eclipseKeywords}, % style of strings
    stepnumber=1,
    numbersep=8pt,
    tabsize=2,
    showstringspaces=false,
    breaklines=true,
    string=[s]{"}{"},
    comment=[l]{:\ "},
    morecomment=[l]{:"},
    literate=
        *{0}{{{\color{numb}0}}}{1}
         {1}{{{\color{numb}1}}}{1}
         {2}{{{\color{numb}2}}}{1}
         {3}{{{\color{numb}3}}}{1}
         {4}{{{\color{numb}4}}}{1}
         {5}{{{\color{numb}5}}}{1}
         {6}{{{\color{numb}6}}}{1}
         {7}{{{\color{numb}7}}}{1}
         {8}{{{\color{numb}8}}}{1}
         {9}{{{\color{numb}9}}}{1}
}


%% Preamble
  \title{Cross-lingual Training for\\Multiple-Choice Question Answering}
  \author{\textbf{Guillermo Echegoyen Blanco} \\ \'Alvaro Rodrigo \\ Anselmo Pe\~nas \\ 
    \{gblanco, alvarory, anselmo\} \textbf{at} lsi.uned.es \\ \\
    NLP \& IR Group \\
    National Distance Education University (UNED)}
    \titlegraphic{\hfill\includegraphics[height=0.75cm]{./figures/UNED.jpeg}
  }
  \date{}
  \addtobeamertemplate{frametitle}{}{%
    \begin{tikzpicture}[remember picture,overlay]
      \node[anchor=north east,yshift=-8pt,xshift=-8pt] at (current page.north east) {{\scriptsize$|$ UNED}};
    \end{tikzpicture}
  }


\begin{document}

\maketitle

\begin{frame}{Overview}
  \setbeamertemplate{section in toc}[sections numbered]
  \tableofcontents[hideallsubsections]
\end{frame}

\section{Introduction}
  \begin{frame}{\secname}
    \begin{alertblock}{Multiple-Choice Question Answering}
      \mAlertSpace
      \textbf{Def:} Given a supporting text, a question and a set of possible answers, choose the correct one. Commonly used to measure reading comprehension in humans.
    \end{alertblock}
    \begin{itemize}
      \item The majority of datasets are in English.
      \item Non-English datasets are usually small.
      \item Usually extracted from exams for humans.
    \end{itemize}
    \alert{ToDo := Add example?, name collections?}
  \end{frame}

  \begin{frame}{\secname}
    \begin{alertblock}{Motivation}
      \begin{itemize}
        \item How to zero-shot transfer from a big MC-QA collection to a smaller one.
        \item Can we zero-shot transfer to a smaller collection in another language?
        \item Harder exams for humans are so for machines too?
      \end{itemize}
    \end{alertblock}
  \end{frame}

\section{Problem Statement}
  \begin{frame}{\secname}
    \begin{alertblock}{Datasets}
      \begin{itemize}
        \item \textbf{RACE} (\cite{Lai2017}): Collected from Chinese schools English exams. \textbf{$>$ 97K questions} with, 4 possible answers each, English monolingual.
        \item \textbf{Entrance Exams} (\cite{rodrigo_systems_2018}): University access exams in Japan. \textbf{$\approx$ 200 questions}, 4 possible answers each. Crowd-translated to 4 different languages.
      \end{itemize}
    \end{alertblock}
    \metroset{block=fill}
    \begin{alertblock}{Example (taken from RACE)}
      \textbf{Evidence:} "The park is open from 8 am to 5 pm." \\
      \textbf{Question:} The park is open for \_\_ hours a day. \\
      \textbf{Options:} A. eight B. nine C. ten D. eleven
    \end{alertblock}
  \end{frame}

  \note[itemize]{%
    \item RACE:
      \begin{enumerate}
        \item Divided into two collections: middle and high school.
      \end{enumerate}
    \item EE:
      \begin{enumerate}
        \item $\approx$ 500 times smaller than RACE. 
        \item Not suitable for fine tuning.
        \item Translated to Spanish, Italian, French, Russian, and (just one edition) German.
        \item Exams from three different years
      \end{enumerate}
    \item RACE are pre University exams, EE are exams at University level.
  }

  \begin{frame}{\secname}
    \begin{alertblock}{Models \& Baselines}
      \mAlertSpace
      \begin{columns}[T, onlytextwidth]
        \column{0.35\textwidth}
          \mAlertSpace
          \begin{itemize}
            \item BERT-base
            \item Multi BERT-base
          \end{itemize}
        \column{0.65\textwidth}
          \mAlertSpace
          \begin{itemize}
            \item Random
            \item Longest answer (\cite{Rogers2020quail})
          \end{itemize}
      \end{columns}
    \end{alertblock}
    \mAlertSpace
    \begin{alertblock}{Method}
      \begin{itemize}
        \item No hyper-parameters search.
        \item Fine-tune each model over RACE.
        \item Test each model over RACE.
        \item Test each model over Entrance Exams in all languages and all years
      \end{itemize}
    \end{alertblock}
  \end{frame}

\section{Results}
  \begin{frame}{\secname}
    % Results table
{\centering
\begin{tabular}{l|c|c|c|c}
  \hline
  Dataset             & BERT     & MultiBERT        & Random   & Longest  \\ \hline
  RACE Mid            &  0.5265  & \textbf{0.6114}  & 0.2500   & 0.3078   \\
  RACE High           &  0.4774  & \textbf{0.5031}  & 0.2500   & 0.3059   \\
  RACE All            &  0.4917  & \textbf{0.5347}  & 0.2500   & 0.3059   \\ \hline
  EE English          &  0.4921  & \textbf{0.4974}  & 0.2500   & 0.2304   \\
  EE Spanish          &  0.3665  & \textbf{0.4503}  & 0.2500   & 0.2932   \\
  EE Italian          &  0.2880  & \textbf{0.4293}  & 0.2500   & 0.2775   \\
  EE French           &  0.3037  & \textbf{0.4346}  & 0.2500   & 0.2565   \\
  EE Russian          &  0.2618  & \textbf{0.3403}  & 0.2500   & 0.2723   \\
  EE German**           &  0.3708  & \textbf{0.4494}  & 0.2500   & 0.2584   \\ \hline
\end{tabular}
}

    \mAlertSpace
    \small{** German only available for one year.}
  \end{frame}

  \note[itemize]{%
    \item Multi-BERT performs better in all scenarios
    \item Performance drops with human-graded difficulty (Mid $<$ High $<$ EE)
    \item Russian is specially difficult as language with very different semantics are not well understood nor tokenized by the model.
  }

\section{Conclusions \& Future Work}
  \begin{frame}{\secname}
    \begin{alertblock}{Conclusions}
      \begin{itemize}
        \item Zero-shot transfer to a smaller task still holds performance in the same language.
        \item Can be done to a different task and language with a multilingual model.
        \item Performance is hampered by exams difficulty in the same way human grades do.
      \end{itemize}
    \end{alertblock}
    \begin{alertblock}{Future Work}
      \begin{itemize}
        \item Continue exploring low-resource languages.
      \end{itemize}
      \alert{ToDo:= Remove Future work? Add something else?}
    \end{alertblock}
  \end{frame}

\section{Outcomes}
  \begin{frame}{\secname}
    \begin{alertblock}{Out main contributions are:}
      \begin{itemize}
        \item SOTA on Entrance Exams in several languages.
        \item RACE trained BERT and Multi BERT models.
      \end{itemize}
    \end{alertblock}
  \end{frame}

  \note[itemize]{%
    \item SOTA in Spanish, Italian and German. 
  }

\begin{frame}[standout]
  Thank you!\\
  Questions?
\end{frame}

\section{References}
  \begin{frame}[allowframebreaks]{References}
    \printbibliography%
  \end{frame}

\end{document}
